% \begin{document}
\chapter{Funzioni Simmetriche}

\section{Serie formali di Potenze}
\dfn{Anello in piu' variabili}{TODO: era tipo $ \mathbb{Q}[x_1,x_2,...,x_n] $}
\dfn{Anello delle serie formali}{
  L'anello delle serie formali su $ \mathbb{Q} $ (va bene qualsiasi anello TODO: commutativo?) e'
  \[
    \mathbb{Q}[[X]] = \{\sum_{n\geq 0} a_nx^n | a_n \in \mathbb{Q}\} 
  \]

Somma e prodotto sono definiti in modo classico:
$$\left(\sum_n a_n x^n\right) + \left(\sum_n b_n x^n\right) = \sum_n (a_n + b_n) x^n$$
$$\left(\sum_n a_n x^n\right) \cdot \left(\sum_n b_n x^n\right) = \sum_n \left(\sum_{k=0}^n a_k b_{n-k}\right) x^n$$
}
Gli elementi si possono pensare come successioni $(a_0, a_1, a_2, \ldots)$ di coefficienti razionali, o equivalentemente come funzioni $\mathbb{N} \to \mathbb{Q}$.

\nt{
Attenzione: sono "polinomi infiniti". Non ci interessa la convergenza — tutto è formale. La parola "formale" significa che lavoriamo puramente con le regole algebriche, senza analisi.
}

\mprop{Invertibilita di serie formali}{
  $f = \sum_{n \geq 0} a_n x^n$ è invertibile in $\mathbb{Q}[[x]]$ se e solo se $a_0 \neq 0$.
}
\ex{Serie invertibile}{
  La serie $1 + ax + a^2x^2 + \cdots = \sum_{n \geq 0} a^n x^n$ è l'inverso di $(1-ax)$, perché:

$$(1-ax)\left(\sum_{n \geq 0} a^n x^n\right) = 1$$

Per questo scriviamo $\dfrac{1}{1-ax}$ come notazione abbreviata per quella serie — non è una frazione nel senso analitico, è solo un simbolo comodo.
}

Possiamo definire anche una derivata.
\dfn{Derivata}{
  $ D: \mathbb{Q}[[x]] \to \mathbb{Q}[[x]] $ e' definita come:
  $$D\left(\sum_{n \geq 0} a_n x^n\right) := \sum_{n \geq 1} n\, a_n\, x^{n-1}$$
}
Valgono le regole usuali, dimostrabili per puro calcolo algebrico (senza limiti):
\begin{itemize}
\item $D(f + g) = Df + Dg$
\item $D(fg) = Df \cdot g + f \cdot Dg$ 
\end{itemize}
Dimostrazioni lasciate a ddestroyer04.

Possiamo definire anche esponenziale e logaritmo:
\dfn{Esponenziale}{
Se $ f(x) \in \mathbb{Q}[[x]] $ allora
\[
  e^{f(x)} = exp(f(x)) \coloneq \sum_{n \geq 0} \frac{1}{n!}f(x)^n \in \mathbb{Q}[[x]]
\]
}

\dfn{Logaritmo}{
Per $g \in 1 + x\,\mathbb{Q}[[x]]$ (cioè con termine costante 1):
$$\log(g(x)) := \sum_{n \geq 1} (-1)^{n-1} \frac{(g(x)-1)^n}{n}$$
In particolare $\log(1+u) = u - \dfrac{u^2}{2} + \dfrac{u^3}{3} - \cdots$ (la serie familiare, ma formale).
}

\mprop{Proprieta' derivata e esponenziale/logaritmo}{
Valgono quello che ci aspettiamo, non perche' stiamo tirando in gioco l'analisi ma puramente per le definizioni date.
$$\log(\exp(f)) = f$$
$$\exp(\log(1+xf)) = 1+xf$$
$$\exp(f+g) = \exp(f)\cdot\exp(g)$$
$$\log(g \cdot h) = \log(g) + \log(h) \quad \text{(quando definiti)}$$
$$D(\exp(f(x))) = f'(x) \cdot \exp(f(x))$$
$$D(\log(g(x))) = \frac{D(g(x))}{g(x)}$$  
}

\section{Polinomi Simmetrici}
Torniamo ora a $ \mathbb{Q}[x_1,...,x_k] $ (anello di polinomi in $ n $ indeterminate). $ S_k $ agisce su $ \mathbb{Q}[x_1,...,x_k] $:
\[
  \sigma f(x_1,...,x_k) = f(x_{\sigma_1},...,x_{\sigma_k})
\]
permutando le variabili.

\dfn{Anello delle funzioni simmetriche}{
  L'anello delle funzioni simmetriche in $k$ variabili è:
  \[
    Sym[X_k] \coloneq \mathbb{Q}[x_1,...,x_k]^{S_k}\ \text{(punti fissi sotto l'azione di )}S_k
  \]
  cioè i polinomi invarianti sotto qualsiasi permutazione delle variabili.
}

\ex{Polinomi Simmetrici}{$x_1^2 + x_2^2 + x_3^2$ è simmetrico. $x_1^2 x_2 + x_2^2 x_3 + x_3^2 x_1$ non è simmetrico (si verifica applicando la trasposizione $(1\,3)$).}

\mprop{Decomposizione per grado}{
  Poiché $S_k$ preserva il grado dei monomi:

$$\mathrm{Sym}[X_k] = \bigoplus_{n \geq 0} \mathrm{Sym}^n[X_k]$$

dove $\mathrm{Sym}^n[X_k]$ sono i polinomi simmetrici omogenei di grado $n$.
}

Notiamo che l'azione di $ S_k $ rispetta il grado.

Supponiamo che $ k \geq n $. No aspe' sta roba qua' serve per sopra, i $ k $ dopo non centrano.

Una composizione di $ n $ 
\dfn{Composizione}{
  \[
    \alpha = (\alpha_1,...,\alpha_n)\ \alpha_i \in \mathbb{N} \setminus \{0\}
  \]
  tali che
  \[
  \sum_{i = 1}^{k} \alpha_i = n
  \]
  Cosa cambia dalle weak compositions? Le weak possono essere 0

  $ \alpha_i $ e' una \textit{parte} di $ \alpha $ e una composiziione con $ k $ parti si dice \textit{k-composizione}.
}

\ex{Composizioni}{
  Le composizioni di 4 sono:
}

La funzione dalle weak k-compositions di $ n $ alle compositions $ n + k $:
\[
  (\alpha_1,...,\alpha_k) \mapsto (\alpha_1 + 1,..., \alpha_k + 1)
\]
e' una biezione. Quindi la cardinalita' di questi insiemi e' uguale, ed e' (ce lo dice dopo perche' si e' dimenticata)

E' $ \binom{n + k -1}{k-1} $, la dimostrazione segue dal lemma:
\mlenma{Dimensione compositions}{
  $ \{\text{k-composizioni di } n\} $ e' in corrispondenza biunivoca con $ \{\text{sottoinsiemi di } \{1,..., n-1\} \} $ di cardinalita' $ k - 1 $:
  \[
    (\alpha_1,...,\alpha_k) \gets \to \{\alpha_1, \alpha_1 + \alpha_2,...,\alpha_1+\alpha_2+...+\alpha_{k-1}\}
  \]
}
boh fine non so, qualcosa ha cucinato. Ora andiamo a vedere le funzioni simmetriche che ci servono di piu'.

\section{Funzioni simmetriche}
Sia $ k \geq n $, definiremo 4 basi diverse di $ Sym^n[X_k] $.

Sia $ \lambda \vdash n $ e sia $ \alpha $ una weak composition di $ n $.

Allora ad $ \alpha $ possiamo associare una partizione di $ n $. E' abbastanza ovvio, non formalizziamo tutto, vediamo un esempio:
\ex{}{
  $ \alpha = (3,1,0,1,2) $ e' una comp (debole) di 7, allora $ \lambda(\alpha) = (3,2,1,1) $. Basta togliere gli 0 e permutare dai e' ez.
}

Sia $ \lambda \vdash n $ data.

Sia $ m_\lambda(x) = \sum_{\alpha \models n} x^\alpha $ ($ \alpha $ e' una k-composizione debole di $ n $), dove $ x^\alpha = x_1^{\alpha_1}...x_k^{\alpha_k} $.

\ex{}{
  $ \lambda=(2,1,1) $, $ n = 4 $ e $ k = 3 $, allora
  \[
    m_{(2,1,1)} = x_1^2x_2x_3 + x_1x_2^2x_3 + x_1x_2x_3^2
  \]
  Chiaramente questo e' un polinomio simmetrico: $ m_\lambda \in Sym^n[X_k] $.
}
La marta sta avendo un esaurimento nervoso, ci fara' sapere perche'. Comunque questo e'.

Possiamo scrivere
\[
  f = \sum_{\alpha \models n} c_\alpha x^\alpha \in Sym^n[X_k]
\]
e' un generico elemento dell'insieme, quindi e' solo un altro modo per descriverlo.

A ogni composizione di $ n $ corrisponde un monomio, a qui corrisponde una partizione.

Vabbe' la marta si e' messa a piangere e ha cancellato tutto, le dispiace. Ricomincia:

se $ f $ e' un polinomio in $ x_1,...,x_k $ omogeneo di grado $ n $, allora $ f $ e' somma di monomi i cui esponenti delle variabili formano una weak composition, a cui possiamo associare una partizione. Ha ridetto la stessa cosa, ok.

Se $ f $ e' anche simmetrico ($ f \in Sym^n[X_k] $) allora $ f = \sum c_\lambda m_\lambda  $.
\ex{}{
  $ n = 2 $ e $ k = 3 $, $ f = 3x_1x_2 + 3x_2x_3 + 3 x_1x_3 =  $ monomi con stessa $ \lambda $ avranno anche lo stesso coefficente quindi possiamo raggrupparli.
}

Gli $ m_\lambda $ sono lin. ind., quindi $ dim Sym^n[X_k] = \# Par(n) $ e generano, quindi sono una base.

Se $ A = (a_{ij})_{i,j \geq 1} $ con un numero finito di coefficenti non nulli. $ r_i = \sum_{j} a_{ij} $ (somma elementi riga $ i $) e $ c_j = \sum_{i} a_{ij} $ (somma elementi colonna $ j $).

$ row A = (r_1,r_2,...) $ e $ col A = (c_1,c_2,...) $. Una $ 0-1 $ matrice e' una matrice a coeff. in $ \{0,1\} $.

\dfn{}{
  $ e_s = m_{1^s} = \sum_{i_1 < i_2 < ... < i_s} x_{i_1}x_{i_2}...x_{i_s} $
}
\dfn{}{
  Sia $ \lambda \vdash n $, $ e_\lambda = e_{\lambda_1}e_{\lambda_2}...e_{\lambda_k} $
}

\ex{}{
  Se $ \lambda = (2,2,1) $ ($ k = 3 $):

  \[
    e_1 = m_{(1,1,1)} = x_1 + x_2 + x_3
  \]
  \[
    e_2 = (x_1x_2+x_1x_3+x_2x_3)
  \]
  \[
  e_{\lambda} = e_2^2 + e_1
  \]
}

\mprop{}{
  Sia $ \lambda \vdash n $ e sia $ \alpha = (\alpha_1,...) $ una composizione debole di $ n $.
  \[
    e_{\lambda} \in Sym^n[X_k] \implies e_\lambda = \sum_{\mu \vdash n} M_{\lambda \mu}m_\mu
  \]
  Dove $ M_{\lambda \mu} $ e' il numero di $ (0,1) $ matrici con $ row(A) = \lambda $ e $ col(A) = \alpha $ e $ \mu = \mu(\alpha) $
}


\section{Funzioni di potenza $ p_\lambda $}

\dfn{Funzioni simmetriche di potenza}{
  Per ogni $ n \geq 1 $:
  \[
    p_n := \sum_i x_i^n = m_{(n)} .
  \]
  Per convenzione $ p_0 = 1 $. Data $ \lambda = (\lambda_1, \dots, \lambda_k) \vdash n $:
  \[
    p_\lambda := p_{\lambda_1}\cdot p_{\lambda_2} \cdots p_{\lambda_k} .
  \]
}

\ex{$ \lambda = (2,1) $, $ k = 3 $}{
  \begin{align*}
    p_{(2,1)} &= (x_1^2 + x_2^2 + x_3^2)(x_1 + x_2 + x_3) \\
    &= x_1^3 + x_2^3 + x_3^3 + x_1^2 x_2 + x_1^2 x_3 + x_1 x_2^2 + x_2^2 x_3 + x_1 x_3^2 + x_2 x_3^2 .
  \end{align*}
}

\subsection*{Funzione generatrice}

Definiamo
\[
  P(t) := \sum_{n \geq 1} \frac{p_n}{n}\, t^n .
\]
La scelta del fattore $ \frac{1}{n} $ rende le formule con $ \exp $ e $ \log $ pulite. Da $ H(t) = \prod_i \frac{1}{1 - x_i t} $:
\[
  \log H(t) = \sum_i \log \frac{1}{1 - x_i t} = \sum_i \sum_{n \geq 1} \frac{x_i^n t^n}{n} = \sum_{n \geq 1} \frac{p_n}{n} t^n = P(t) ,
\]
da cui
\[
  \boxed{\ H(t) = \exp(P(t))\ } .
\]
Analogamente, da $ E(t) = \prod_i (1 + x_i t) $:
\[
  \log E(t) = \sum_i \sum_{n \geq 1} \frac{(-1)^{n - 1}}{n} x_i^n t^n = -P(-t) ,
\]
da cui
\[
  \boxed{\ E(t) = \exp(-P(-t))\ } .
\]

\section{Identita' di Newton}

\thm{Identita' di Newton (omogenee)}{
  Per ogni $ m \geq 1 $:
  \[
    m\, h_m = \sum_{i = 1}^m p_i\, h_{m - i} .
  \]
}
\pf{Dim}{
  Da $ H(t) = \exp(P(t)) $, derivando rispetto a $ t $:
  \[
    H'(t) = P'(t) \cdot \exp(P(t)) = P'(t) \cdot H(t) .
  \]
  Esplicitamente
  \[
    H'(t) = \sum_{n \geq 1} n\, h_n\, t^{n - 1}, \qquad P'(t) = \sum_{n \geq 1} p_n\, t^{n - 1} .
  \]
  L'identita' $ H'(t) = P'(t) H(t) $ diventa
  \[
    \sum_{n \geq 1} n h_n t^{n - 1} = \Big( \sum_{n \geq 1} p_n t^{n - 1} \Big) \Big( \sum_{n \geq 0} h_n t^n \Big) .
  \]
  Estraendo il coefficiente di $ t^{m - 1} $: $ m h_m = \sum_{i = 1}^m p_i h_{m - i} $.
}

\thm{Identita' di Newton (elementari)}{
  Per ogni $ m \geq 1 $:
  \[
    m\, e_m = \sum_{i = 1}^m (-1)^{i - 1}\, p_i\, e_{m - i} .
  \]
}
\pf{Dim}{
  Deriviamo $ E(t) = \exp(-P(-t)) $: $ E'(t) = P(-t) \cdot E(t) $. Notare il segno: $ \frac{d}{dt}(-P(-t)) = P'(-t) $, e $ P(-t) = \sum_{n \geq 1} (-1)^{n - 1} \frac{p_n}{n} t^n $... estraendo il coefficiente di $ t^{m - 1} $ si ottiene la formula.
}

\ex{Calcolo di $ h_2 $ e $ e_2 $}{
  Per $ m = 2 $: $ 2 h_2 = p_1 h_1 + p_2 h_0 = p_1^2 + p_2 $, quindi
  \[
    h_2 = \frac{p_1^2 + p_2}{2} .
  \]
  Analogamente $ 2 e_2 = p_1 e_1 - p_2 e_0 = p_1^2 - p_2 $, quindi
  \[
    e_2 = \frac{p_1^2 - p_2}{2} .
  \]
  Le due differiscono solo per il segno di $ p_2 $.
}

\section{Le $ p_\lambda $ sono una base}

\cor{}{
  $ \{p_\lambda \mid \lambda \vdash n,\ \ell(\lambda') \leq k\} $ e' una base di $ Sym^n[X_k] $.
}
\pf{Dim (traccia)}{
  Si ordinano gli $ e_\lambda $ con l'ordine lessicografico inverso. Le identita' di Newton esprimono $ p_n $ in termini di $ e_1, \dots, e_n $, con termine dominante $ (-1)^{n - 1} n\, e_n $. Quindi per ogni $ \lambda $ il termine dominante di $ p_\lambda $ e' proporzionale a $ e_\lambda $ con coefficiente non nullo: la matrice di transizione da $ \{p_\lambda\} $ a $ \{e_\lambda\} $ e' triangolare con diagonale non nulla, quindi invertibile. Poiche' $ \{e_\lambda\} $ e' una base, anche $ \{p_\lambda\} $ lo e'.
}

\subsection*{Sviluppo di $ h_m $ nelle $ p_\lambda $}

Da $ H(t) = \exp(P(t)) $ si ricava
\[
  h_m = \sum_{\lambda \vdash m} \frac{1}{z_\lambda} p_\lambda
\]
dove $ z_\lambda := \prod_{i \geq 1} i^{\alpha_i}\, \alpha_i! $, con $ \alpha_i = \#\{j : \lambda_j = i\} $.

\osservazione{Interpretazione combinatoria.
  In $ S_m $ il numero di permutazioni di tipo di ciclo $ \lambda $ e' $ \frac{m!}{z_\lambda} $. Quindi
  \[
    h_m = \frac{1}{m!} \sum_{\sigma \in S_m} p_{\lambda(\sigma)} .
  \]
  E' una media sul tipo di ciclo di tutte le permutazioni di $ S_m $.
}

\ex{$ m = 2 $}{
  Le partizioni di 2 sono $ (2) $ e $ (1,1) $. $ z_{(2)} = 2^1 \cdot 1! = 2 $, $ z_{(1,1)} = 1^2 \cdot 2! = 2 $. Quindi
  \[
    h_2 = \frac{p_2}{2} + \frac{p_1^2}{2} = \frac{p_1^2 + p_2}{2} \checkmark
  \]
}

\subsection*{Effetto di $ \omega $ sulle $ p_n $}

\mprop{}{
  $ \omega(p_n) = (-1)^{n - 1} p_n $.
}
\pf{Dim}{
  Applichiamo $ \omega $ alle identita' di Newton per $ h_m $:
  \[
    m\, \omega(h_m) = \sum_{i = 1}^m \omega(p_i)\, \omega(h_{m - i}) .
  \]
  Poiche' $ \omega(h_m) = e_m $:
  \[
    m\, e_m = \sum_{i = 1}^m \omega(p_i)\, e_{m - i} .
  \]
  Confrontando con l'identita' di Newton per $ e_m $: $ m\, e_m = \sum_{i = 1}^m (-1)^{i - 1}\, p_i\, e_{m - i} $. Le due coincidono se e solo se $ \omega(p_i) = (-1)^{i - 1} p_i $.
}

Per $ \lambda \vdash n $ con $ \ell(\lambda) = k $:
\[
  \omega(p_\lambda) = \prod_{j = 1}^k (-1)^{\lambda_j - 1} p_{\lambda_j} = (-1)^{\sum_j (\lambda_j - 1)} p_\lambda = (-1)^{n - \ell(\lambda)} p_\lambda .
\]

\osservazione{Schema riassuntivo (relazioni fra basi).
  \[
    H(t) = \exp(P(t)), \qquad E(t) = \exp(-P(-t)), \qquad H(t) \cdot E(-t) = 1
  \]
  \[
    h_m = \sum_{\lambda \vdash m} \frac{p_\lambda}{z_\lambda}, \qquad m h_m = \sum_i p_i h_{m - i}, \qquad m e_m = \sum_i (-1)^{i - 1} p_i e_{m - i}
  \]
  \[
    \omega(e_n) = h_n, \quad \omega(h_n) = e_n, \quad \omega(p_n) = (-1)^{n - 1} p_n, \quad \omega^2 = id .
  \]
}

\section{Raccordo: le quattro basi e il limite $ Sym[X] $}

\subsection*{Ricapitolo delle basi}

Abbiamo definito tre famiglie di generatori per $ \mathbb{Q}[x_1, \dots, x_k]^{S_k} $:
\[
  e_m = \sum_{i_1 < \dots < i_m} x_{i_1} \cdots x_{i_m}, \quad
  h_m = \sum_{i_1 \leq \dots \leq i_m} x_{i_1} \cdots x_{i_m}, \quad
  p_m = \sum_i x_i^m
\]
e quattro basi per $ Sym^n[X_k] $:
\[
  m_\lambda = \sum_{\lambda(\alpha) = \lambda} x^\alpha, \qquad
  e_\lambda = \prod e_{\lambda_i}, \qquad
  h_\lambda = \prod h_{\lambda_i}, \qquad
  p_\lambda = \prod p_{\lambda_i} .
\]

\osservazione{
  Le relazioni algebriche tra queste funzioni \textbf{non dipendono dal numero di variabili} $ k $ (purche' $ k $ sia abbastanza grande). Ad esempio
  \[
    e_2 = \frac{p_1^2 - p_2}{2} .
  \]
  Verifica per $ k = 2 $: $ e_2 = x_1 x_2 = \frac{(x_1 + x_2)^2 - x_1^2 - x_2^2}{2} $. Verifica per $ k = 3 $: $ e_2 = x_1 x_2 + x_1 x_3 + x_2 x_3 = \frac{(x_1 + x_2 + x_3)^2 - x_1^2 - x_2^2 - x_3^2}{2} $. La formula e' la stessa per ogni $ k \geq 2 $.
}

\osservazione{Coefficienti interi.
  \begin{itemize}
    \item Le $ m_\lambda $ generano $ \mathbb{N}[x_1, \dots, x_k]^{S_k} $ con coefficienti in $ \mathbb{N} $.
    \item Le $ e_\lambda $ e $ h_\lambda $ generano $ \mathbb{Z}[x_1, \dots, x_k]^{S_k} $ con coefficienti in $ \mathbb{Z} $.
    \item Le $ p_\lambda $ generano $ \mathbb{Z}[x_1, \dots, x_k]^{S_k} $ con coefficienti in $ \mathbb{Z} $.
  \end{itemize}
}

\subsection*{Il limite: $ Sym[X] $}

Fissata $ \lambda \vdash n $, la funzione $ m_\lambda(x_1, \dots, x_k) $ dipende da $ k $, ma al crescere di $ k $ si ottiene una catena di mappe compatibili
\[
  \cdots \to Sym[X_k] \to Sym[X_{k-1}] \to \cdots
\]
dove ogni freccia e' la specializzazione $ x_k = 0 $. Il \textbf{limite inverso} di questa catena e' l'anello delle \textit{funzioni simmetriche in infinite variabili}:
\[
  Sym[X] := \varprojlim_k Sym[X_k] .
\]
Possiamo lavorare con un numero infinito di variabili $ x_1, x_2, \dots $, e le basi $ \{m_\lambda\}, \{e_\lambda\}, \{h_\lambda\}, \{p_\lambda\}, \{s_\lambda\} $ restano basi anche in questo contesto.

\section{La funzione di Schur come quoziente di determinanti}

\subsection*{Polinomi alternanti}

\dfn{Polinomio alternante}{
  Un polinomio $ f \in \mathbb{Q}[x_1, \dots, x_k] $ e' \textit{alternante} se per ogni $ \sigma \in S_k $:
  \[
    \sigma \cdot f = \mathrm{sgn}(\sigma)\, f .
  \]
}

\textbf{Costruzione.} Data una sequenza $ \lambda = (\lambda_1, \dots, \lambda_k) $ di interi, definiamo
\[
  a_\lambda := \det(x_i^{\lambda_j})_{i,j = 1}^k = \sum_{\sigma \in S_k} \mathrm{sgn}(\sigma)\, x_1^{\lambda_{\sigma(1)}} \cdots x_k^{\lambda_{\sigma(k)}} .
\]
E' alternante per costruzione (scambiare due righe del determinante cambia il segno). Se due componenti di $ \lambda $ coincidono, il determinante ha due colonne uguali e si annulla; quindi $ a_\lambda \neq 0 $ solo se i $ \lambda_i $ sono tutti distinti.

\subsection*{Il determinante di Vandermonde}

Il caso fondamentale e' $ \delta = (k-1, k-2, \dots, 1, 0) $:
\[
  a_\delta = \det(x_i^{k - j})_{i,j = 1}^k = \prod_{1 \leq i < j \leq k} (x_i - x_j) .
\]

\ex{$ k = 3 $}{
  \[
    a_\delta = (x_1 - x_2)(x_1 - x_3)(x_2 - x_3) = \det\begin{pmatrix} x_1^2 & x_1 & 1 \\ x_2^2 & x_2 & 1 \\ x_3^2 & x_3 & 1 \end{pmatrix} .
  \]
}

\mprop{Caratterizzazione dei polinomi alternanti}{
  Un polinomio $ p $ e' alternante se e solo se e' divisibile per $ a_\delta $.
}
\pf{Dim}{
  $ (\Leftarrow) $ Se $ p = a_\delta \cdot g $ con $ g $ simmetrico, allora $ \sigma(p) = \mathrm{sgn}(\sigma)\, a_\delta \cdot g = \mathrm{sgn}(\sigma)\, p $.

  $ (\Rightarrow) $ Se $ p $ e' alternante, per ogni coppia $ i \neq j $ la trasposizione $ (ij) $ da' $ (ij) \cdot p = -p $. Specializzando $ x_i = x_j $: $ p|_{x_i = x_j} = -p|_{x_i = x_j} $, quindi $ p|_{x_i = x_j} = 0 $, cioe' $ (x_i - x_j) \mid p $. Poiche' questi fattori sono a due a due coprimi e il loro prodotto e' $ a_\delta $, si ha $ a_\delta \mid p $.
}

\cor{}{
  $ f $ e' alternante $ \iff f = a_\delta \cdot g $ con $ g $ simmetrico. L'applicazione $ g \mapsto a_\delta \cdot g $ e' un \textit{isomorfismo} $ Sym[X_k] \xrightarrow{\sim} A_k $ fra simmetrici e alternanti.
}

\subsection*{Definizione algebrica della funzione di Schur}

\dfn{Funzione di Schur (algebrica)}{
  Per ogni $ \lambda \vdash n $ con $ \ell(\lambda) \leq k $, la \textit{funzione di Schur} e'
  \[
    s_\lambda := \frac{a_{\lambda + \delta}}{a_\delta} ,
  \]
  dove $ \lambda + \delta = (\lambda_1 + k - 1,\ \lambda_2 + k - 2,\ \dots,\ \lambda_k + 0) $.
}

E' un polinomio simmetrico perche' numeratore e denominatore sono entrambi alternanti (corollario sopra).

\ex{$ s_{(2,1)} $ con $ k = 2 $}{
  $ \lambda = (2,1) $, $ \delta = (1,0) $, $ \lambda + \delta = (3,1) $:
  \[
    a_{(3,1)} = \det\begin{pmatrix} x_1^3 & x_1 \\ x_2^3 & x_2 \end{pmatrix} = x_1^3 x_2 - x_1 x_2^3 = x_1 x_2 (x_1 - x_2)(x_1 + x_2)
  \]
  \[
    a_{(1,0)} = x_1 - x_2
  \]
  \[
    s_{(2,1)} = \frac{x_1 x_2 (x_1 - x_2)(x_1 + x_2)}{x_1 - x_2} = x_1 x_2 (x_1 + x_2) = x_1^2 x_2 + x_1 x_2^2 = m_{(2,1)} .
  \]
  Con $ k = 3 $ si verifica $ s_{(2,1)} = m_{(2,1)} + 2 m_{(1,1,1)} $.
}

\thm{$ \{s_\lambda\} $ e' una base}{
  $ \{s_\lambda \mid \lambda \vdash n,\ \ell(\lambda) \leq k\} $ e' una base di $ Sym^n[X_k] $.
}
\pf{Dim}{
  \textbf{Passo 1.} Mostriamo che $ \{a_{\lambda + \delta} \mid \lambda \vdash n,\ \ell(\lambda) \leq k\} $ e' una base dello spazio $ A_k^{(m)} $ degli alternanti omogenei di grado $ m = n + |\delta| $.

  Sia $ f = \sum_\alpha c_\alpha x^\alpha \in A_k^{(m)} $ alternante. Se $ c_\alpha \neq 0 $, allora $ \alpha_i \neq \alpha_j $ per $ i \neq j $ (altrimenti la trasposizione $ (ij) $ darebbe $ f = -f $, quindi $ f = 0 $). Gli esponenti distinti si ordinano come $ \alpha = \lambda + \delta $ per una partizione $ \lambda $. L'alternanza forza $ c_{\sigma(\alpha)} = \mathrm{sgn}(\sigma)\, c_\alpha $, quindi
  \[
    f = \sum_\lambda c_{\lambda + \delta}\, a_{\lambda + \delta} .
  \]

  \textbf{Passo 2.} La mappa $ \phi: Sym[X_k] \to A_k $, $ g \mapsto a_\delta \cdot g $, e' un isomorfismo che manda $ s_\lambda \mapsto a_\delta \cdot s_\lambda = a_{\lambda + \delta} $.

  \textbf{Passo 3.} Poiche' $ \phi $ e' un isomorfismo e $ \{a_{\lambda + \delta}\} $ e' una base di $ A_k^{(m)} $, la preimmagine $ \{s_\lambda\} $ e' una base di $ Sym^n[X_k] $.
}

\cor{Estrazione dei coefficienti di Schur}{
  Se $ f \in Sym[X_k] $ si scrive $ f = \sum_\mu c_\mu\, s_\mu $, allora
  \[
    c_\lambda = [x^{\lambda + \delta}]\ a_\delta \cdot f ,
  \]
  cioe' il coefficiente di $ s_\lambda $ in $ f $ e' il coefficiente del monomio $ x^{\lambda + \delta} $ nel prodotto $ a_\delta \cdot f $.
}
\pf{Dim}{
  Moltiplicando $ f = \sum_\mu c_\mu\, s_\mu $ per $ a_\delta $:
  \[
    a_\delta \cdot f = \sum_\mu c_\mu\, a_\delta\, s_\mu = \sum_\mu c_\mu\, a_{\mu + \delta} .
  \]
  Il monomio $ x^{\lambda + \delta} $ appare solo in $ a_{\lambda + \delta} $ (con coefficiente 1, essendo il monomio ``di testa'' del determinante), quindi il suo coefficiente in $ a_\delta \cdot f $ e' esattamente $ c_\lambda $.
}

\section{Definizione combinatoria di Schur tramite SSYT}

\dfn{Peso di un tableau e Schur combinatoria}{
  Dato un SSYT $ T $ di forma $ \lambda $ con valori in $ \{1, \dots, k\} $, il \textit{peso} e'
  \[
    x^T := \prod_{i = 1}^k x_i^{\#\{i \in T\}} .
  \]
  La \textit{funzione di Schur combinatoria} e'
  \[
    \widehat{s}_\lambda := \sum_{T \in SSYT(\lambda),\ \text{valori in } \{1, \dots, k\}} x^T .
  \]
}

\ex{$ \lambda = (2,1) $, $ k = 3 $}{
  I SSYT di forma $ (2,1) $ con valori in $ \{1,2,3\} $ sono 8:
  \[
    \begin{ytableau} 2 \\ 1 & 1 \end{ytableau} \quad
    \begin{ytableau} 3 \\ 1 & 1 \end{ytableau} \quad
    \begin{ytableau} 2 \\ 1 & 2 \end{ytableau} \quad
    \begin{ytableau} 3 \\ 1 & 2 \end{ytableau} \quad
    \begin{ytableau} 2 \\ 1 & 3 \end{ytableau} \quad
    \begin{ytableau} 3 \\ 1 & 3 \end{ytableau} \quad
    \begin{ytableau} 3 \\ 2 & 2 \end{ytableau} \quad
    \begin{ytableau} 3 \\ 2 & 3 \end{ytableau}
  \]
  con monomi $ x_1^2 x_2,\ x_1^2 x_3,\ x_1 x_2^2,\ x_1 x_2 x_3,\ x_1 x_2 x_3,\ x_1 x_3^2,\ x_2^2 x_3,\ x_2 x_3^2 $.
  Sommando:
  \[
    \widehat{s}_{(2,1)} = x_1^2 x_2 + x_1^2 x_3 + x_1 x_2^2 + x_2^2 x_3 + x_1 x_3^2 + x_2 x_3^2 + 2 x_1 x_2 x_3 = m_{(2,1)} + 2 m_{(1,1,1)} = s_{(2,1)} \checkmark
  \]
}

Le due definizioni coincidono: e' il \textit{Teorema fondamentale delle funzioni di Schur} (verra' dimostrato piu' avanti come $ s_\lambda = \widehat{s}_\lambda $).

\section{Corrispondenza RSK: introduzione e algoritmo}

\subsection*{Motivazione}

La corrispondenza RSK (Robinson--Schensted--Knuth) e' una biezione fondamentale in combinatoria. Risponde alla domanda: c'e' un modo canonico per associare tableau a matrici di interi non negativi?

\thm{RSK}{
  Esiste una biezione
  \[
    \{\mathbb{N}\text{-matrici finite}\} \xleftrightarrow{\;\sim\;} \bigsqcup_\lambda \{\text{coppie } (R, L) \text{ di SSYT di forma } \lambda\} .
  \]
}

\subsection*{Parola di lettura di una matrice}

Data una matrice $ A = (a_{ij}) $ a coefficienti in $ \mathbb{N} $, costruiamo la \textit{parola di lettura} $ w_A $: per ogni colonna $ j $ da sinistra a destra, scriviamo il valore di riga $ i $ ripetuto $ a_{ij} $ volte, scorrendo le righe dall'alto in basso.

\ex{}{
  \[
    A = \begin{pmatrix} 1 & 0 & 2 \\ 0 & 2 & 0 \\ 1 & 0 & 0 \end{pmatrix}
  \]
  Colonna 1: riga 1 da' 1 una volta, riga 3 da' 3 una volta $ \to 1, 3 $.
  Colonna 2: riga 2 da' 2 due volte $ \to 2, 2 $.
  Colonna 3: riga 1 da' 1 due volte $ \to 1, 1 $.
  \[
    w_A = 1, 3, 2, 2, 1, 1 .
  \]
}

\subsection*{Algoritmo di inserzione: il bumping}

L'ingrediente fondamentale e' il \textit{bumping} (scacciamento): come inserire un numero $ t $ in una riga di un SSYT preservando la struttura.

\textbf{Algoritmo di inserzione di $ t $ nella riga $ r $:}
\begin{enumerate}
  \item Se $ t $ e' maggiore o uguale di tutti gli elementi della riga, lo aggiungo alla fine. Fine.
  \item Altrimenti, trovo il primo elemento $ s > t $ nella riga. Sostituisco $ s $ con $ t $: $ t $ ``scaccia'' $ s $.
  \item Inserisco $ s $ nella riga successiva con lo stesso algoritmo.
  \item Quando un elemento si aggiunge alla fine di una riga, si aggiunge anche una nuova cella al tableau.
\end{enumerate}

\textbf{Proprieta' chiave del bumping.}
\begin{itemize}
  \item Se inserisco $ k $ in un SSYT $ T $, ottengo ancora un SSYT.
  \item Il percorso di inserzione va verso il basso: ogni elemento scacciato finisce nella riga successiva piu' in basso.
  \item L'elemento scacciato e' sempre strettamente maggiore di quello che lo scaccia, quindi nelle righe successive l'elemento inserito e' piu' grande, e puo' scacciare qualcosa piu' a sinistra.
\end{itemize}

\ex{Bumping}{
  Inseriamo $ 4 $ nel tableau
  \[
    \begin{ytableau} 8 & 9 \\ 4 & 4 & 6 & 7 \\ 2 & 3 & 3 & 5 & 6 & 8 \\ 1 & 1 & 2 & 4 & 5 & 6 \end{ytableau}
  \]
  $ 4 $ in riga 1: primo elemento $ > 4 $ e' il secondo $ 5 $. Sostituisco $ 5 $ con $ 4 $.

  $ 5 $ in riga 2: primo elemento $ > 5 $ e' $ 6 $. Sostituisco $ 6 $ con $ 5 $.

  $ 6 $ in riga 3: primo elemento $ > 6 $ e' $ 7 $. Sostituisco $ 7 $ con $ 6 $.

  $ 7 $ in riga 4: $ 7 < 8 $ e $ 7 < 9 $... no, $ 7 > 8 $? No. $ 7 $ si aggiunge alla fine della riga 4 dopo $ 9 $? In realta', si applica la regola: $ 7 $ si confronta col primo elemento della riga 4 ($ 8 $); $ 8 > 7 $, quindi $ 7 $ scaccia $ 8 $. $ 8 $ va in riga 5.
}

\subsection*{L'algoritmo RSK}

\textbf{Input.} La parola $ w_A = w_1, w_2, \dots, w_N $.

\textbf{Processo.} Manteniamo due tableau $ R $ (di \textit{inserzione}) e $ L $ (di \textit{registrazione}), entrambi inizialmente vuoti. Per ogni $ w_j $:
\begin{enumerate}
  \item Inserisco $ w_j $ in $ R $ tramite bumping.
  \item In $ L $, aggiungo il numero di colonna $ j $ (cioe' l'indice della parola, ovvero il valore corrispondente in $ w_A $) nella nuova cella creata da $ R $.
\end{enumerate}

\textbf{Output.} La coppia $ (R, L) $ di SSYT della stessa forma $ \lambda $.

\ex{Esempio completo}{
  $ A = \begin{pmatrix} 1 & 0 & 2 \\ 0 & 2 & 0 \\ 1 & 0 & 0 \end{pmatrix} $, $ w_A = 1, 3, 2, 2, 1, 1 $.

  \textbf{Inserisco $ 1 $:}
  \[
    R = \begin{ytableau} 1 \end{ytableau}, \quad L = \begin{ytableau} 1 \end{ytableau}
  \]

  \textbf{Inserisco $ 3 $:} $ 3 > 1 $, si aggiunge alla fine di riga 1.
  \[
    R = \begin{ytableau} 1 & 3 \end{ytableau}, \quad L = \begin{ytableau} 1 & 1 \end{ytableau}
  \]

  \textbf{Inserisco $ 2 $:} in riga 1 il primo $ > 2 $ e' $ 3 $; lo sostituisco e $ 3 $ va in riga 2.
  \[
    R = \begin{ytableau} 3 \\ 1 & 2 \end{ytableau}, \quad L = \begin{ytableau} 1 \\ 1 & 1 \end{ytableau}
  \]

  \textbf{Inserisco $ 2 $:} attenzione, il bumping usa $ > $ stretto, quindi $ 2 $ non scaccia $ 2 $; si aggiunge in fondo a riga 1.
  \[
    R = \begin{ytableau} 3 \\ 1 & 2 & 2 \end{ytableau}, \quad L = \begin{ytableau} 1 \\ 1 & 1 & 2 \end{ytableau}
  \]

  \textbf{Inserisco $ 1 $:} in riga 1 il primo $ > 1 $ e' il primo $ 2 $; lo sostituisco e $ 2 $ va in riga 2. In riga 2 il primo $ > 2 $ e' $ 3 $; lo sostituisco e $ 3 $ va in riga 3.
  \[
    R = \begin{ytableau} 3 \\ 2 \\ 1 & 1 & 2 \end{ytableau}, \quad L = \begin{ytableau} 2 \\ 1 \\ 1 & 1 & 2 \end{ytableau}
  \]

  \textbf{Inserisco $ 1 $:} in riga 1 $ 1 \not> 1 $, quindi si aggiunge in fondo.
  \[
    R = \begin{ytableau} 3 \\ 2 \\ 1 & 1 & 1 & 2 \end{ytableau}, \quad L = \begin{ytableau} 2 \\ 1 \\ 1 & 1 & 2 & 3 \end{ytableau}
  \]

  Risultato: $ (R, L) $ e' una coppia di SSYT di forma $ \lambda = (4, 1, 1) $.
}

\subsection*{Osservazioni}

\osservazione{
  La corrispondenza RSK e' una biezione: da $ (R, L) $ si puo' ricostruire univocamente $ A $. L'algoritmo inverso esiste ed e' simmetrico.
}

\osservazione{
  Lungo il percorso di inserzione ci si muove sempre verso il basso, perche' l'elemento scacciato e' strettamente maggiore di quello che lo ha scacciato.
}

\osservazione{
  Se inserisco $ k $ in un SSYT $ T $, ottengo ancora un SSYT: l'algoritmo mantiene la struttura ad ogni passo.
}

\subsection*{Perche' RSK e' importante}

\textbf{Collegamento con i polinomi di Schur.} $ s_\lambda $ conta i SSYT di forma $ \lambda $. RSK spiega perche' il prodotto
\[
  s_\mu \cdot s_\nu = \sum_\lambda c^\lambda_{\mu \nu}\, s_\lambda
\]
ha coefficienti $ c^\lambda_{\mu \nu} $ (i numeri di \textit{Littlewood--Richardson}) interi non negativi: contano coppie di tableau.

\textbf{Caso speciale.} Se $ A $ e' una matrice di permutazione, RSK si riduce alla classica corrispondenza di Robinson--Schensted fra permutazioni e coppie di SYT.

\section{RSK: biunivocita' e prima ricostruzione}

\subsection*{Esercizio 8 -- Ricostruzione di $A$ data $(L,R)$}

Dati $ (L,R) $ coppia di SSYT della \textit{stessa forma}, vogliamo ricostruire la $ \mathbb{N} $-matrice finita $ A $ tale che $ \mathrm{RSK}(A) = (L,R) $.

\ex{Caso esplicito}{
  Siano
  \[
    R = \begin{ytableau} 2 & 3 \\ 1 & 1 & 3 \end{ytableau}
    \qquad
    L = \begin{ytableau} 2 & 3 \\ 1 & 2 & 3 \end{ytableau}
    \qquad \text{(notazione francese: riga lunga in basso)}
  \]
  Vogliamo capire qual'e' \textit{l'ultima casella aggiunta}. Procediamo con il bumping inverso, ragionando sulla casella con il valore $ T $ massimo che sta piu' a destra nella riga piu' bassa (regola minima): in questo modo otteniamo passo dopo passo
  \[
    \binom{3}{3},\ \binom{3}{2},\ \binom{2}{1},\ \binom{1}{3},\ \binom{1}{2}
  \]
  e quindi
  \[
    w_A = \begin{pmatrix} 1 & 1 & 2 & 3 & 3 \\ 2 & 3 & 1 & 2 & 3 \end{pmatrix}
    \implies
    A = \begin{pmatrix} 0 & 1 & 1 \\ 1 & 0 & 0 \\ 0 & 1 & 1 \end{pmatrix}
  \]
}

\textbf{Procedura generale per invertire RSK:}
\begin{enumerate}
  \item Trovo le caselle di $ R $ con valore $ T $ massimo.
  \item Quella aggiunta per ultima e' la \textit{piu' a destra nella riga piu' bassa} (cioe' minima per l'ordine di bumping).
  \item Sapendo qual e' la casella aggiunta, leggo in $ L $ il numero che vi era stato inserito; il bumping inverso mi dice quale numero era stato scacciato dalla riga sopra, e cosi' via.
\end{enumerate}

Avevamo gia' dimostrato (osservazioni 1--5) che RSK e' un'applicazione
\[
  \{\mathbb{N}\text{-matrici finite}\} \longrightarrow \{(L,R) \text{ SSYT della stessa forma}\}.
\]

\mprop{RSK e' biunivoca}{
  Dato $ (L,R) $ SSYT della stessa forma, e' sempre possibile ricostruire $ w_A $ e quindi $ A $: ad ogni passo la casella aggiunta e' univocamente determinata dalla regola sopra.
}

\section{Forma bilineare su $ Sym[X] $}

\subsection*{Ripasso: matrici $ \mathbb{N} $}

Avevamo gia' dimostrato che
\[
  e_\lambda = \sum_{\mu \vdash n} M_{\lambda,\mu} m_\mu
\]
dove $ M_{\lambda,\mu} $ e' il coefficiente di $ x^\mu $ in $ e_\lambda $, ovvero il numero di $ (0,1) $-matrici con $ row(A) = \lambda $ e $ col(A) = \mu $.

Si fa lo stesso ragionamento con le $ h_\lambda $: ora pero' in ogni colonna possiamo scegliere variabili \textit{con ripetizione}.

\mprop{Sviluppo di $ h_\lambda $ nelle $ m_\mu $}{
  \[
    h_\lambda = \sum_{\mu \vdash n} N_{\lambda,\mu} m_\mu
  \]
  dove $ N_{\lambda,\mu} $ e' il coefficiente di $ x^\mu $ in $ h_\lambda $, ovvero il numero di $ \mathbb{N} $-matrici con $ row(A) = \lambda $ e $ col(A) = \mu $.
}
\pf{Dim}{
  Si fa nel modo solito: in ogni colonna $ j $ scelgo un multi-insieme di taglia $ \lambda_j $ di variabili, ovvero un vettore $ (a_{1j}, a_{2j}, \dots) \in \mathbb{N}^\infty $ con $ \sum_i a_{ij} = \lambda_j $. Rimpiazzo le variabili con gli esponenti, ottenendo una $ \mathbb{N} $-matrice.
}

\osservazione{
  $ N_{\lambda,\mu} = N_{\mu,\lambda} $: $ A \leftrightarrow A^T $ e' una biezione che scambia $ \lambda $ e $ \mu $. Lo stesso vale per $ M $.
}

\subsection*{Definizione della forma bilineare}

\dfn{Forma bilineare $ \langle \cdot, \cdot \rangle $ su $ Sym[X] $}{
  $ \langle \cdot, \cdot \rangle $ e' l'unica forma bilineare tale che
  \[
    \langle m_\lambda, h_\mu \rangle = \delta_{\lambda,\mu} = \begin{cases} 1 & \text{se } \lambda = \mu \\ 0 & \text{altrimenti} \end{cases}
  \]
  Dato che $ \{m_\lambda\} $ e $ \{h_\mu\} $ sono basi, la richiesta le rende \textit{basi duali}. Ogni forma bilineare puo' essere descritta in questo modo (su un prodotto di basi).
}

\section{Identita' di Cauchy}

\mprop{Identita' di Cauchy}{
  \[
    \prod_{i,j} \frac{1}{1 - x_i y_j} = \sum_{\lambda} m_\lambda(x) h_\lambda(y)
  \]
  e
  \[
    \prod_{i,j} (1 + x_i y_j) = \sum_{\lambda} m_\lambda(x) e_\lambda(y)
  \]
  (la somma e' su tutte le partizioni di tutti gli $ n \in \mathbb{N} $.)
}
\pf{Dim}{
  Dimostriamo la prima identita' (la seconda e' analoga). Sviluppando la serie geometrica:
  \[
    \prod_{i,j} \frac{1}{1 - x_i y_j} = \prod_{i,j}(1 + z_{ij} + z_{ij}^2 + \dots), \qquad z_{ij} = x_i y_j .
  \]
  Un monomio $ x^\alpha y^\beta $ si ottiene scegliendo una $ \mathbb{N} $-matrice $ A = (a_{ij}) $ tale che $ \prod_{i,j}(x_i y_j)^{a_{ij}} = x^\alpha y^\beta $, ovvero con $ row(A) = \alpha $ e $ col(A) = \beta $. Quindi
  \[
    \prod_{i,j} \frac{1}{1-x_i y_j} = \sum_\lambda \sum_\mu N_{\lambda,\mu}\, m_\lambda(x)\, m_\mu(y) = \sum_\lambda m_\lambda(x) \underbrace{\Big(\sum_\mu N_{\lambda,\mu} m_\mu(y) \Big)}_{= h_\lambda(y)} .
  \]
}

\mprop{Caratterizzazione della dualita'}{
  Se $ \{u_\lambda\} $ e $ \{v_\lambda\} $ sono due basi di $ Sym[X] $, allora sono \textit{duali} se e solo se
  \[
    \prod_{i,j} \frac{1}{1 - x_i y_j} = \sum_\lambda u_\lambda(x) v_\lambda(y).
  \]
}
\pf{Dim}{Lasciata al lettore (Macdonald cap. I).}

\mprop{Simmetria della forma bilineare}{
  $ \langle \cdot, \cdot \rangle $ e' simmetrica.
}
\pf{Dim}{
  \[
    \langle h_\lambda, h_\mu \rangle = \Big\langle \sum_\nu N_{\lambda,\nu} m_\nu,\, h_\mu \Big\rangle = \sum_\nu N_{\lambda,\nu} \langle m_\nu, h_\mu \rangle = N_{\lambda,\mu} = N_{\mu,\lambda} = \langle h_\mu, h_\lambda \rangle .
  \]
}

\section{Le $ p_\lambda $ e il prodotto scalare di Hall}

\subsection*{Ripasso: $ z_\lambda $}

Avevamo definito
\[
  z_\lambda = \prod_i (m_i!\, \cdot\, i^{m_i}), \qquad m_i = \#\{j : \lambda_j = i\}.
\]

\ex{$ \lambda = (2,1,1) $}{
  $ m_1 = 2,\ m_2 = 1 $, quindi $ z_\lambda = 2!\cdot 1!\cdot 1^2 \cdot 2^1 = 4 $.
}

E avevamo dimostrato che
\[
  h_n = \sum_{\lambda \vdash n} \frac{1}{z_\lambda} p_\lambda .
\]

\mprop{Identita' di Cauchy nelle $ p_\lambda $}{
  \[
    \prod_{i,j} \frac{1}{1 - x_i y_j} = \sum_\lambda \frac{p_\lambda(x)\, p_\lambda(y)}{z_\lambda}
  \]
  \[
    \prod_{i,j} (1 + x_i y_j) = \sum_\lambda (-1)^{|\lambda| - \ell(\lambda)} \frac{p_\lambda(x)\, p_\lambda(y)}{z_\lambda}
  \]
  Quindi $ \{p_\lambda\} $ e $ \{p_\lambda / z_\lambda\} $ sono basi duali.
}
\pf{Dim}{
  Sia $ z_{ij} = x_i y_j $. Con la serie geometrica solita
  \[
    \prod_{i,j} \frac{1}{1 - x_i y_j t} = \prod_{i,j}(1 + z_{ij} t + z_{ij}^2 t^2 + \dots) = \sum_{n \in \mathbb{N}} h_n(z) t^n = \sum_n \Big( \sum_{\lambda \vdash n} \frac{1}{z_\lambda} p_\lambda(z) \Big) t^n
  \]
  Valutando in $ t = 1 $:
  \[
    \prod_{i,j} \frac{1}{1 - x_i y_j} = \sum_\lambda \frac{p_\lambda(z)}{z_\lambda} .
  \]
  Resta da osservare che, con $ z_{ij} = x_i y_j $, vale per ogni $ n $:
  \[
    p_n(z) = \sum_{i,j} (x_i y_j)^n = \Big( \sum_i x_i^n \Big) \Big( \sum_j y_j^n \Big) = p_n(x) p_n(y) .
  \]
  Quindi $ p_\lambda(z) = p_\lambda(x) p_\lambda(y) $.
}

\cor{Prodotto scalare di Hall}{
  $ \langle \cdot, \cdot \rangle $ e' un \textit{prodotto scalare} (definito positivo) e $ \{p_\lambda\} $ e' una base ortogonale, con
  \[
    \langle p_\lambda, p_\mu \rangle = z_\lambda\, \delta_{\lambda,\mu} .
  \]
}
\pf{Dim}{
  Dalla dualita' di $ \{p_\lambda\} $ e $ \{p_\mu / z_\mu\} $: $ \langle p_\lambda, p_\mu / z_\mu \rangle = \delta_{\lambda,\mu} $, da cui $ \langle p_\lambda, p_\mu \rangle = z_\lambda \delta_{\lambda,\mu} $. E' definito positivo perche' gli $ z_\lambda $ sono positivi.
}

\cor{$ \omega: h_\lambda \leftrightarrow e_\lambda $ e' un'isometria}{}
\pf{Dim}{
  Basta verificare $ \langle \omega(p_\lambda), \omega(p_\mu) \rangle = \langle p_\lambda, p_\mu \rangle $. Sappiamo che $ \omega(p_n) = (-1)^{n-1} p_n $, quindi $ \omega(p_\lambda) = (-1)^{|\lambda|-\ell(\lambda)} p_\lambda $ e
  \[
    \langle \omega(p_\lambda), \omega(p_\mu) \rangle = (\pm 1)^2\, z_\lambda\, \delta_{\lambda,\mu} = \langle p_\lambda, p_\mu \rangle .
  \]
}

\section{Funzioni di Schur combinatorie $ \widehat{s}_\lambda $}

\dfn{Funzioni di Schur combinatorie}{
  \[
    \widehat{s}_\lambda(x) = \sum_{T \in SSYT(\lambda)} x^T
  \]
  dove $ x^T = \prod_i x_i^{\#\{i \in T\}} $. Si estende analogamente a forme skew $ \lambda/\mu $.
}

\thm{Le $ \widehat{s}_\lambda $ sono una base ortonormale di $ Sym[X] $}{}
\pf{Dim}{
  Dobbiamo mostrare che $ \langle \widehat{s}_\lambda, \widehat{s}_\mu \rangle = \delta_{\lambda,\mu} $. Equivalentemente (una base e' duale a se stessa $ \iff $ e' ortonormale): basta far vedere che
  \[
    \sum_\lambda \widehat{s}_\lambda(x)\, \widehat{s}_\lambda(y) = \prod_{i,j} \frac{1}{1 - x_i y_j} .
  \]
  Un monomio $ x^\alpha y^\beta $ nello sviluppo del prodotto di Cauchy si ottiene scegliendo una $ \mathbb{N} $-matrice $ A $ con $ row(A) = \alpha $ e $ col(A) = \beta $; quindi il coefficiente di $ x^\alpha y^\beta $ nel secondo membro e' il numero di tali matrici. Nel primo membro, e' il numero di coppie di SSYT $ (L,R) $ della stessa forma $ \lambda $ tali che $ type(R) = \alpha $ e $ type(L) = \beta $. Per la biezione RSK, i due numeri sono uguali.
}

\section{Numeri di Kostka}

\dfn{Numeri di Kostka}{
  Siano $ \lambda/\mu $ una forma skew (eventualmente con $ \mu = \emptyset $) e $ \alpha \models n $ una composizione. Si pone
  \[
    K_{\lambda/\mu, \alpha} = \#\{T \text{ SSYT di forma } \lambda/\mu \mid type(T) = \alpha\} .
  \]
  Caso particolare $ \alpha = (1^n) $:
  \[
    f^{\lambda/\mu} := K_{\lambda/\mu,\, 1^n} = \#\{\text{tableaux \textit{standard} di forma } \lambda/\mu\} .
  \]
}

\mprop{}{
  \[
    \sum_{\lambda \vdash n} (f^\lambda)^2 = n!
  \]
}
\pf{Dim}{
  Restringiamo la biezione RSK alle \textit{matrici di permutazione}: $ \sigma \in S_n $ corrisponde alla matrice $ A_\sigma $ con $ w_{A_\sigma} = \binom{1\,2\cdots n}{\sigma(1)\,\sigma(2)\cdots\sigma(n)} $. Otteniamo cosi' una biezione
  \[
    S_n \xleftrightarrow{\,RSK\,} \bigsqcup_{\lambda \vdash n} SYT(\lambda) \times SYT(\lambda) ,
  \]
  da cui $ n! = \sum_\lambda (f^\lambda)^2 $.
}

\ex{$ n = 4 $}{
  Le partizioni di 4 con i rispettivi $ f^\lambda $:
  \[
    f^{(4)} = 1, \quad f^{(3,1)} = 3, \quad f^{(2,2)} = 2, \quad f^{(2,1,1)} = 3, \quad f^{(1^4)} = 1
  \]
  e $ 1^2 + 3^2 + 2^2 + 3^2 + 1^2 = 24 = 4! $. (Ricorda: somma dei quadrati delle dimensioni degli irriducibili di $ S_4 $.)
}

\thm{}{
  \[
    e_\mu = \sum_{\lambda \vdash |\mu|} K_{\lambda'/\mu}\, s_\lambda
  \]
  (la somma usa $ \lambda' $, la trasposta). La dimostrazione non fa parte del programma.
}

\cor{Regola di Pieri}{
  \[
    s_\nu \cdot e_n = \sum_\lambda s_\lambda
  \]
  dove la somma e' sulle partizioni $ \lambda $ tali che $ \lambda/\nu $ e' una \textit{striscia verticale} (cioe' $ n $ caselle, mai due nella stessa riga).
}

\ex{$ \nu = (3,3,1) $, $ n = 2 $}{
  Dobbiamo trovare tutte le posizioni in cui aggiungere 2 caselle, mai due nella stessa riga. Si ottiene
  \[
    s_\nu \cdot e_2 = s_{(4,4,1)} + s_{(3,3,1,1,1)} + s_{(4,3,2)} + s_{(4,3,1,1)} + s_{(3,3,2,1)} .
  \]
}

\section{Simmetria e proprieta' di $ \widehat{s}_\lambda $}

\mprop{Le $ \widehat{s}_\lambda $ sono simmetriche}{}
\pf{Dim}{
  Le trasposizioni adiacenti $ (i, i+1) $ generano $ S_\infty $, quindi basta mostrare $ (i, i+1) \cdot \widehat{s}_{\lambda/\mu} = \widehat{s}_{\lambda/\mu} $ (l'azione scambia $ i $ con $ i+1 $).

  Costruiamo una biezione fra SSYT di forma $ \lambda/\mu $ che scambia $ \#\{i\} $ e $ \#\{i+1\} $:
  \begin{itemize}
    \item Le colonne che non contengono ne' $ i $ ne' $ i+1 $ non cambiano.
    \item Le colonne che li contengono entrambi sono gia' incastrate verticalmente (con $ i $ sopra a $ i+1 $) e non si toccano.
    \item Restano le colonne con \textit{solo uno} tra $ i $ e $ i+1 $: in ogni riga esse formano blocchi del tipo
    \[
      \underbrace{i\ i\ \cdots\ i}_{r}\ \underbrace{i{+}1\ \cdots\ i{+}1}_{s}
    \]
    e li sostituiamo con $ s $ entrate $ i $ seguite da $ r $ entrate $ i+1 $. La condizione di SSYT e' preservata.
  \end{itemize}
}

\cor{}{
  \[
    s_{\lambda/\mu}(x) = \sum_{\nu \vdash |\lambda/\mu|} K_{\lambda/\mu, \nu}\, m_\nu(x)
  \]
  In particolare, per $ \mu = \emptyset $:
  \[
    \widehat{s}_\lambda(x) = \sum_{\nu \vdash |\lambda|} K_{\lambda, \nu}\, m_\nu(x) .
  \]
}
\pf{Dim}{
  Per definizione $ \widehat{s}_{\lambda/\mu} = \sum_\alpha K_{\lambda/\mu,\alpha}\, x^\alpha $ con $ \alpha $ composizione. Per la simmetria appena dimostrata, le composizioni che differiscono per permutazione hanno lo stesso coefficiente $ K_{\lambda/\mu, \nu} $ (con $ \nu $ partizione associata), quindi raggruppando si ottengono le $ m_\nu $.
}

\subsection*{Ordine di dominanza}

Ricordiamo che $ \mu \trianglelefteq \lambda $ se per ogni $ i \geq 1 $ vale $ \sum_{j=1}^i \mu_j \leq \sum_{j=1}^i \lambda_j $.

\mprop{}{
  Siano $ \lambda, \mu \vdash n $. Se $ K_{\lambda,\mu} \neq 0 $ allora $ \mu \trianglelefteq \lambda $. Inoltre $ K_{\lambda,\lambda} = 1 $.
}
\pf{Dim}{
  Sia $ T $ un SSYT di forma $ \lambda $ e tipo $ \mu $. In un SSYT tutti gli $ 1 $ stanno sulla prima riga (sotto violerebbero la stretta crescenza in colonna), quindi
  \[
    \mu_1 = \#\{1 \in T\} \leq \lambda_1 .
  \]
  Allo stesso modo gli $ 1 $ e i $ 2 $ devono comparire solo nelle prime due righe, quindi
  \[
    \mu_1 + \mu_2 \leq \lambda_1 + \lambda_2 .
  \]
  Cosi' via, $ \mu_1 + \dots + \mu_i \leq \lambda_1 + \dots + \lambda_i $ per ogni $ i $.

  Per $ K_{\lambda,\lambda} = 1 $: l'unico SSYT di forma $ \lambda $ e tipo $ \lambda $ e' quello in cui la riga $ i $ contiene solo $ i $ (super-standard).
}

\cor{Le $ \widehat{s}_\lambda $ sono una base di $ Sym[X] $}{}
\pf{Dim}{
  La matrice di cambio di base $ K = (K_{\lambda,\mu}) $ e' triangolare (rispetto all'ordine di dominanza) con $ 1 $ sulla diagonale, quindi invertibile.
}

\osservazione{
  Riassumendo: le $ \widehat{s}_\lambda $ con $ \lambda \vdash n $ sono una base \textit{ortonormale} di $ Sym^n[X] $.
}

\section{Sviluppi nelle $ \widehat{s}_\lambda $}

\cor{}{
  \[
    h_\mu = \sum_\lambda K_{\lambda, \mu}\, \widehat{s}_\lambda
  \]
}
\pf{Dim}{
  Sia $ h_\mu = \sum_\lambda a_{\lambda,\mu}\, \widehat{s}_\lambda $. Per ortonormalita' delle $ \widehat{s}_\lambda $ e dualita' $ \langle h_\mu, m_\nu \rangle = \delta_{\mu,\nu} $:
  \[
    a_{\lambda,\mu} = \langle h_\mu, \widehat{s}_\lambda \rangle = \Big\langle h_\mu, \sum_\nu K_{\lambda,\nu} m_\nu \Big\rangle = \sum_\nu K_{\lambda,\nu}\, \delta_{\mu,\nu} = K_{\lambda,\mu} .
  \]
}

\cor{}{
  \[
    h_1^n = \sum_{\lambda \vdash n} f^\lambda\, \widehat{s}_\lambda .
  \]
}
\pf{Dim}{
  $ h_1^n = h_{(1^n)} $ e applicando il corollario precedente con $ \mu = (1^n) $: $ K_{\lambda, 1^n} = f^\lambda $.
}

\section{Trasposta e simmetria di RSK}

\osservazione{(Dimostrazione lasciata a Morbidelli)
  Sia $ A $ una $ \mathbb{N} $-matrice finita e $ A \xrightarrow{RSK} (L,R) $. Allora $ A^T \xrightarrow{RSK} (R, L) $.
}

\cor{}{
  $ A $ e' simmetrica $ \iff R = L $. Inoltre RSK ristretta e' una biezione fra $ \mathbb{N} $-matrici simmetriche con $ row(A) = \alpha $ e SSYT di tipo $ \alpha $.
}

\cor{Involuzioni di $ S_n $}{
  \[
    \sum_{\lambda \vdash n} f^\lambda = |\{\sigma \in S_n \mid \sigma^2 = id\}|
  \]
}
\pf{Dim}{
  Le matrici di permutazione sono ortogonali, quindi $ A^{-1} = A^T $, ossia $ \sigma = \sigma^{-1} $ se e solo se $ \sigma^2 = id $. Restringendo RSK alle matrici di permutazione \textit{simmetriche} otteniamo la biezione fra involuzioni di $ S_n $ e SYT di forma $ \lambda \vdash n $.
}

\ex{$ n = 4 $}{
  $ \sum_\lambda f^\lambda = 1 + 3 + 2 + 3 + 1 = 10 $. Le involuzioni di $ S_4 $: $ 6 $ trasposizioni + $ 3 $ doppie trasposizioni + $ 1 \cdot id = 10 $. Torna.
}

\section{RSK duale (RSK*)}

\textbf{Algoritmo.} Biezione fra $ (0,1) $-matrici finite e coppie $ (L,R) $ di tableaux della stessa forma tali che $ L' $ ed $ R $ sono SSYT (cioe' $ L $ e' SSYT trasposto: strettamente crescente in riga, debolmente crescente in colonna). Funziona nello stesso modo solo che, durante l'inserimento di un numero $ j $, si scaccia l'elemento piu' a sinistra che e' $ \leq j $.

Cosi' come abbiamo dimostrato (con RSK)
\[
  \prod_{i,j} \frac{1}{1 - x_i y_j} = \sum_\lambda \widehat{s}_\lambda(x)\, \widehat{s}_\lambda(y) ,
\]
allo stesso modo, con RSK*, si dimostra che
\[
  \prod_{i,j} (1 + x_i y_j) = \sum_\lambda \widehat{s}_\lambda(x)\, \widehat{s}_{\lambda'}(y) .
\]

\cor{Azione di $ \omega $ su $ \widehat{s}_\lambda $}{
  \[
    \omega(\widehat{s}_\lambda) = \widehat{s}_{\lambda'}
  \]
  (analogamente per le $ \widehat{s}_{\lambda/\mu} $.)
}
\pf{Dim}{
  Indichiamo con $ \omega_y $ l'involuzione che agisce solo sulle variabili $ y $ (le $ x $ diventano costanti). Calcoliamo:
  \begin{align*}
    \sum_\lambda \widehat{s}_\lambda(x)\, \omega_y(\widehat{s}_\lambda(y)) &= \omega_y\Big(\sum_\lambda \widehat{s}_\lambda(x)\, \widehat{s}_\lambda(y)\Big) \stackrel{Cauchy}{=} \omega_y\Big( \prod_{i,j} \frac{1}{1 - x_i y_j} \Big) \\
    &= \omega_y \Big( \sum_\lambda m_\lambda(x)\, h_\lambda(y) \Big) = \sum_\lambda m_\lambda(x)\, e_\lambda(y) \\
    &= \prod_{i,j} (1 + x_i y_j) = \sum_\lambda \widehat{s}_\lambda(x)\, \widehat{s}_{\lambda'}(y) .
  \end{align*}
  Dato che vale grado per grado e in ogni grado le $ \widehat{s}_\lambda $ formano una base, i coefficienti coincidono: $ \omega(\widehat{s}_\lambda) = \widehat{s}_{\lambda'} $.
}

\section{Identificazione $ s_\lambda = \widehat{s}_\lambda $}

\thm{}{
  \[
    s_\lambda = \widehat{s}_\lambda \qquad \forall\, \lambda .
  \]
  Cioe' la funzione di Schur \textit{algebrica} (definita ad esempio come quoziente di determinanti alla Jacobi-Trudi, o tramite la formula bialternante) coincide con quella \textit{combinatoria}.
}
\pf{Dim}{
  Dal teorema sui Kostka:
  \[
    e_\mu = \sum_\lambda K_{\lambda', \mu}\, s_\lambda .
  \]
  Dall'analoga relazione combinatoria $ h_\mu = \sum_\lambda K_{\lambda, \mu}\, \widehat{s}_\lambda $, applicando $ \omega $:
  \[
    e_\mu = \omega(h_\mu) = \omega\Big( \sum_\lambda K_{\lambda, \mu} \widehat{s}_\lambda \Big) = \sum_\lambda K_{\lambda, \mu}\, \widehat{s}_{\lambda'} = \sum_{\lambda'} K_{\lambda', \mu}\, \widehat{s}_\lambda
  \]
  (riordinando: $ \lambda $ e $ \lambda' $ percorrono entrambi le partizioni di $ |\mu| $.)

  Confrontando le due espressioni per $ e_\mu $:
  \[
    \sum_\lambda K_{\lambda', \mu}\, s_\lambda = \sum_\lambda K_{\lambda', \mu}\, \widehat{s}_\lambda .
  \]
  La matrice $ K $ e' triangolare con $ 1 $ sulla diagonale, quindi invertibile: il sistema ha un'unica soluzione, ovvero $ s_\lambda = \widehat{s}_\lambda $ per ogni $ \lambda $.
}


% \end{document}
